\documentclass[11pt,a4paper]{article}
\usepackage[margin=1.02in]{geometry}
\usepackage[T1]{fontenc}
\usepackage[utf8]{inputenc}
\usepackage{lmodern,microtype}
\usepackage{amsmath,amssymb,amsthm,mathtools,mathrsfs}
\usepackage{booktabs,array,enumitem}
\usepackage[hidelinks]{hyperref}
\usepackage{enumitem}
\allowdisplaybreaks[2]
\setlength{\parskip}{0.35em}
\setlength{\parindent}{0pt}
\newtheorem{theorem}{Theorem}[section]
\newtheorem{proposition}[theorem]{Proposition}
\newtheorem{lemma}[theorem]{Lemma}
\newtheorem{corollary}[theorem]{Corollary}
\theoremstyle{definition}
\newtheorem{definition}[theorem]{Definition}
\newtheorem{remark}[theorem]{Remark}
\newcommand{\FF}{\mathbb F_2}
\newcommand{\QQ}{\mathbb Q}
\newcommand{\ZZ}{\mathbb Z}
\newcommand{\rank}{\operatorname{rank}}
\newcommand{\corank}{\operatorname{corank}}
\newcommand{\Sel}{\operatorname{Sel}_2}
\newcommand{\Sha}{\operatorname{\mathrlap{\mathchoice{\raisebox{.13ex}{$\displaystyle\amalg$}}{\raisebox{.13ex}{$\textstyle\amalg$}}{\amalg}{\amalg}}\amalg}}
\newcommand{\sgnbit}[2]{\left[\frac{#1}{#2}\right]}
\newcommand{\one}{\boldsymbol{1}}
\newcommand{\diag}{\operatorname{diag}}
\newcommand{\Pic}{\operatorname{Pic}}
\title{A row-5a audit and a simultaneous-twist strengthening\\
\large On Manvir Jaswal's congruent-number route to Diophantine definability}
\author{Independent mathematical audit and new deduction}
\date{13 November 2026}
\begin{document}
\maketitle

\begin{abstract}
We examine the arithmetic contribution in M.~Jaswal,
\emph{A proof that $\ZZ$ is Diophantine in $\QQ$ through congruent-number twists},
hexagon:2611.00331v2 (12 November 2026).  The change to congruent-number curves,
minimal $2$-Selmer groups, R\'edei matrices and Tian--Yuan--Zhang's theorem already
belong to the Glazer--McCarthy project.  Jaswal's different \emph{prescribed
Selmer class} ($\beta\equiv1\pmod8$), its first-genus-sum (row-5a) determinant,
and the associated auxiliary-prime pattern constitute a substantive alternative
not found as a completed construction in the two inspected repositories.
We give a proof-level audit of this contribution and prove a strengthening:
\emph{one and the same squarefree twist parameter works for any finite set of
pairwise-coprime admissible auxiliary twists}.  The resulting curves all have
analytic and Mordell--Weil rank exactly one, have minimal $2$-Selmer group,
have trivial $2$-primary Tate--Shafarevich group, and realize the same
prescribed $2$-descent class by infinite-order rational points.  There are
$\gg X/\log X$ suitable parameters up to $X$.  We spell out the Selmer
matrix, finite-field determinant reduction, prime construction, and sources;
we also give an exact numerical configuration as a reproducibility check.
No improvement in the number of variables of a Diophantine polynomial is
asserted.
\end{abstract}

\section{Scope and the novelty question}\label{sec:scope}

The preprint \cite{Jaswal} states the existential Diophantine conclusion
\(\exists N\ge1\;\exists D\in\ZZ[X,Y_1,\dots,Y_N]\), but does \emph{not}
provide an explicit value of $N$, an explicit defining polynomial, a degree
bound, or coefficient bounds.  Its introduction calls Theorem~1.2 its one
new theorem.  Its six-contact height estimate (Theorem~1.3), its prime
configuration theorem (Theorem~1.4), and its elliptic-index formula follow
the earlier Glazer--OpenAI development, albeit with some modifications in
presentation and imported modularity inputs.

The following distinction is essential for correct attribution:
\begin{enumerate}[label=\textup{(\roman*)},itemsep=2pt]
\item The congruent-number family $E_d:y^2=x^3-d^2x$ and a
      R\'edei--Monsky Selmer-minimality construction occur in the
      \href{https://github.com/elliotglazer/diophantine-z-q/blob/main/docs/research/tyz-5b.md}{Glazer--McCarthy row-5b research note}
      and its Lean development.  There the prescribed first square class has
      absolute residue $7\pmod8$, and the decisive genus sum is the
      \emph{second} Tian--Yuan--Zhang sum (row 5b).
\item The \emph{definition} of the \emph{first} genus sum already occurs in
      \href{https://github.com/elliotglazer/diophantine-z-q/blob/main/blueprint/roadmap/rank/literature/genus-sum-5a.md}{the existing row-5a blueprint},
      and both sums occur in the project's work on Proposition~N.
      Therefore the first sum or its source theorem is \emph{not} newly
      discovered by the preprint.
\item What appears new relative to the completed routes inspected is the
      \emph{specific simultaneous existence construction}
      $(\beta,d)\in\Sel(E_e)\cap\Sel(E_{e\theta})$ with
      $\beta\equiv1\pmod8$, certified by row 5a and a compatible
      finite-field configuration.  This is a new combination of known
      ingredients and can replace a harder arithmetic dependency in a
      pole-completeness proof; it is not by itself a new Diophantine
      variable bound.
\end{enumerate}
The comparison includes the main branches of
\href{https://github.com/JoseBrox/diophantine-z-q_improvements}{the user's improvements repository}
and \href{https://github.com/elliotglazer/diophantine-z-q}{Glazer's repository},
plus the open pull-request titles at the time of inspection.  An exhaustive
claim about unmerged private work is not made.

The mathematics proved below is a new strengthening of the preprint's
Theorem~1.2, not a new proof of the entire Diophantine-definability
machinery in its Sections~2--5 and~7--8.

\section{Statement of the strengthened theorem}\label{sec:main}

For $n>0$ squarefree, write
\[
 E_n:\quad y^2=x^3-n^2x=x(x-n)(x+n),\qquad
 \partial_n:E_n(\QQ)/2E_n(\QQ)\hookrightarrow
       (\QQ^\times/\QQ^{\times2})^2,
\]
where, away from $2$-torsion, $\partial_n(x,y)=(x,x-n)$.
Let $\Sigma_{5a}(n)$ denote the \emph{first genus sum} of
Tian--Yuan--Zhang for $n\equiv5\pmod8$ (defined in
\S\ref{sec:published}).

\begin{definition}\label{def:admissible}
For odd positive squarefree $\beta,\theta$, say that $(\beta,\theta)$
is \emph{admissible} if
\[
 \beta>1,\quad \beta\equiv\theta\equiv1\pmod8,\quad
 (\beta,\theta)=1,
\]
and every prime dividing $\beta$ is a quadratic residue of $\theta$
(equivalently $(\theta/p)=+1$ for all $p\mid\beta$), while every
prime dividing $\theta$ is a quadratic residue of $\beta$
(equivalently $(\beta/q)=+1$ for all $q\mid\theta$).
The condition for $\theta=1$ is vacuous.
\end{definition}

\begin{theorem}[Simultaneous prescribed-class twists]\label{thm:sim}
Fix a squarefree $\beta>1$ with $\beta\equiv1\pmod8$, and let
$\theta_1,\ldots,\theta_s$ be any finite family of positive squarefree
integers satisfying:
\begin{enumerate}[label=\textup{(a)},itemsep=2pt]
\item $(\beta,\theta_j)$ is admissible for every $j$;
\item $\gcd(\theta_i,\theta_j)=1$ whenever $i\ne j$.
\end{enumerate}
There are a positive squarefree integer $e$ and a squarefree divisor
$d\mid e/\beta$ with $d\equiv5\pmod8$ such that, setting
$\theta_0=1$ and $n_j=e\theta_j$ for $0\le j\le s$, simultaneously:
\begin{align*}
 &n_j\equiv5\pmod8,\qquad \dim_{\FF}\Sel(E_{n_j})=3,\qquad
          \Sigma_{5a}(n_j)\equiv1\pmod2;\\
 &(\beta,d)\in\Sel(E_{n_j}),\qquad
    \operatorname{ord}_{z=1}L(E_{n_j},z)
      =\rank E_{n_j}(\QQ)=1;\\
 &\Sha(E_{n_j}/\QQ)[2^\infty]=0,\qquad
   \exists P_j\in E_{n_j}(\QQ)\text{ of infinite order with }
     \partial_{n_j}(P_j)=(\beta,d).
\end{align*}
In particular $x(P_j)\in\beta\QQ^{\times2}$ for every $j$.

More precisely, there exist a fixed positive integer $C$, a modulus
$M$ and a reduced residue class $a\pmod M$ such that for every
sufficiently large prime $p\equiv a\pmod M$, the integer $e=Cp$
has all the displayed properties for the same $d$.  Thus the
set of suitable $e$ contains a full prime arithmetic progression
scaled by $C$.

The parameter $e$ can be chosen with exactly
\[
                        \omega(e)=2\omega(\beta)+3
\]
distinct prime factors, independently of $s$, and with all its
auxiliary prime factors arbitrarily large.
For fixed $(\beta;\theta_1,\dots,\theta_s)$ there is a constant
$c>0$ such that
\[
 \#\{e\le X:\ e\text{ has these simultaneous properties}\}
                      \ \ge\ c\frac{X}{\log X}
                  \qquad(X\gg_{\beta,\theta_1,\dots,\theta_s}1).
\]
\end{theorem}

\begin{remark}\label{rem:scope}
The pairwise-coprime assumption makes the free Legendre-symbol bits
for the different $\theta_j$ disjoint.  This proof does not establish
simultaneous realizability for general overlapping families.  It also does
not supply simultaneous rational points with a shared numerical
$y$-coordinate or bound their heights; these are different statements.
\end{remark}

\section{Published analytic input and the genus sum}\label{sec:published}

For positive squarefree $n\equiv5\pmod8$ and a nonempty subset
$T\subseteq\{p:p\mid n\}$ write $n_T=\prod_{p\in T}p$ and
\[
 g(n_T)=\#\bigl(2\operatorname{Cl}(\QQ(\sqrt{-n_T}))\bigr)
                                \pmod2\ \in\FF.
\]
Define
\begin{equation}\label{eq:genus}
 \Sigma_{5a}(n)=
    \sum_{\substack{\pi\text{ partition of }\{p:p\mid n\}\\
           \text{exactly one block }T_0:\ n_{T_0}\equiv5(8)\\
           \text{all other blocks }T:\ n_T\equiv1(8)}}
        \prod_{T\in\pi}g(n_T)\ \in\FF.
\end{equation}
This is Tian--Yuan--Zhang's first sum over unordered factorizations;
the partition with a single block is allowed.  The formula is the
first sum in the published \cite[Theorem~1.4]{TYZ}, not just a
conjectural reading of a preliminary manuscript.

\begin{proposition}[Analytic parity criterion]\label{prop:tyz}
If $n>0$ is squarefree, $n\equiv5\pmod8$, and
$\Sigma_{5a}(n)=1$, then $E_n$ has analytic rank one.
If moreover $\dim_{\FF}\Sel(E_n)=3$, then its Mordell--Weil rank is
one, its $2$-primary Tate--Shafarevich group is zero, and every
class of $\Sel(E_n)$ is represented by a rational point of
infinite order.
\end{proposition}
\begin{proof}
Tian--Yuan--Zhang prove that
$2^{-\rho(n)}\mathcal L(n)$ can be even only if both their
first and second genus sums are even, where $\rho(n)\ge0$ and
$\mathcal L(n)$ vanishes whenever the analytic order exceeds one
\cite[pp.~723--724, Theorem~1.4]{TYZ}.
Consequently $\Sigma_{5a}(n)=1$ implies
$\mathcal L(n)\ne0$, hence analytic order at most one.
The root number is $-1$ for $n\equiv5\pmod8$, so the
analytic order is exactly one.  Gross--Zagier supplies a
nontorsion rational point (the curve is modular, indeed CM),
so $\rank E_n(\QQ)\ge1$.

All three nonidentity points of $E_n[2]$ are rational.  Thus
\[
 \dim_{\FF}E_n(\QQ)/2E_n(\QQ)
        =\rank E_n(\QQ)+2\le\dim_{\FF}\Sel(E_n)=3.
\]
The rank is therefore exactly one, and the injection in the
Kummer sequence
\[
 0\longrightarrow E_n(\QQ)/2E_n(\QQ)
 \overset{\partial_n}{\longrightarrow}\Sel(E_n)
 \longrightarrow \Sha(E_n/\QQ)[2]\longrightarrow0
\]
is surjective.  Hence $\Sha[2]=0$.  Any nonzero element of
$\Sha[2^\infty]$ has a nonzero element of order two, so the
entire $2$-primary torsion subgroup vanishes.
Finally, let a Selmer class be $\partial_n(P)$.
If $P$ is torsion, replace it by $P+2Q$, where $Q$ has
infinite order; the Kummer class does not change and the
new point has infinite order.  This proves the last assertion.
\end{proof}

\section{The exact $2$-Selmer linear system}\label{sec:selmer}

All vector spaces and matrix operations in this section are over
$\FF$.  Write $n=p_1\cdots p_r$, with $n\equiv5\pmod8$,
and set
\[
 y_i=\sgnbit{-1}{p_i},\quad z_i=\sgnbit{2}{p_i},\quad
 a_{ij}=\sgnbit{p_j}{p_i}\quad(i\ne j),\quad
 a_{ii}=\sum_{j\ne i}a_{ij}.
\]
The bracket is the additive Legendre symbol: it equals zero for a
square and one for a nonsquare.  Quadratic reciprocity gives
\begin{equation}\label{eq:reciprocity}
 a_{ij}+a_{ji}=y_i y_j\quad(i\ne j),\qquad
 y\cdot\one=0,\quad z\cdot\one=1.
\end{equation}
Let $A=(a_{ij})$ and $D_y=\diag(y)$, $D_z=\diag(z)$.
The row sums of $A$ vanish, and summing the $j$th column
using reciprocity gives
\begin{equation}\label{eq:Arels}
 A\one=0,\qquad A^\top\one=y,\qquad
 A+A^\top=yy^\top+D_y.
\end{equation}

\begin{lemma}[Monsky--R\'edei system]\label{lem:selmer}
Every element of $\Sel(E_n)$ has a unique representation
\[
 \left((-1)^\varepsilon\prod_i p_i^{u_i},\ 
       (-1)^\varepsilon 2^{z\cdot u}\prod_i p_i^{w_i}\right)
\]
where $(\varepsilon,u,w)\in\FF\times\FF^r\times\FF^r$ satisfies
\begin{align}
 Au+D_y w&=\varepsilon y,                                      &&\tag{S1}\label{eq:S1}\\
 (A+D_y)w+D_z u+(z\cdot u)z&=\varepsilon y,                    &&\tag{S2}\label{eq:S2}\\
 y\cdot w&=y\cdot u.                                             &&\tag{S3}\label{eq:S3}
\end{align}
Set
\begin{equation}\label{eq:M}
 M=\begin{pmatrix}A+A^\top&A^\top\\ A&D_z\end{pmatrix},\qquad
 b=\binom{y+z}{0},\qquad
 M_{5a}=\begin{pmatrix}M&b\\b^\top&0\end{pmatrix}.
\end{equation}
Then
\begin{equation}\label{eq:Selrank}
                \dim_{\FF}\Sel(E_n)=2+\corank_{\FF}M.
\end{equation}
\end{lemma}
\begin{proof}
At $\infty$ the two Kummer coordinates have the same sign.
For odd $p\nmid n$ they are unramified.  For $p\mid n$,
the two local classes $\partial_n(0,0)$ and
$\partial_n(n,0)$ have valuations $(0,1)$ and $(1,0)$;
the unit-square comparisons give precisely the rows of
\eqref{eq:S1} and \eqref{eq:S2}.  At $p=2$, since
$n\equiv5\pmod8$, the curve is locally a quadratic-scaling
of $E_5$.  The classes of $(0,0)$, $(5,0)$ and $(-4,6)$ on
$E_5$ give $(\!-1,-5)$, $(5,2)$ and $(\!-1,-9)$;
they independently generate the local Kummer image of order
$8$.  If $a=2^{v_2(a)}a'$ with $a'$ a $2$-adic unit,
let $s(a)=[a'\equiv3\pmod4]$ and
$t(a)=[a'\equiv\pm3\pmod8]$.  The local image has
$v_2(d_1)=0$, $v_2(d_2)=t(d_1)$, $s(d_2)=s(d_1)$.
This yields the exponent $z\cdot u$, and \eqref{eq:S3}.
The uniqueness is uniqueness of the square-class exponents
supported on $-1,2,p_1,\ldots,p_r$.

For the dimension identity, restrict the solution space to
$\varepsilon=0$, $z\cdot u=1$.  The affine map
\[
 \Xi(u,w)=\bigl(u+w+(1+y\cdot u)\one,\quad u+\one\bigr)
\]
is a bijection onto $\ker M$.  To check one direction,
put $c=1+y\cdot u$, $U=u+w+c\one$, $W=u+\one$.
By \eqref{eq:S1}--\eqref{eq:S2} and \eqref{eq:Arels},
\[
 AU+D_zW
 =Au+Aw+D_zu+z
 =D_yw+(D_yw+D_zu+z)+D_zu+z=0.
\]
Moreover $A^\top w=A w+y(y\cdot w)+D_yw$; thus
\[
 (A+A^\top)U+A^\top W
          =y(y\cdot w+c+1)=y\bigl(y\cdot(u+w)\bigr)=0.
\]
Conversely, the two block equations for $(U,W)\in\ker M$
give $y\cdot U=0$ and $z\cdot W=0$, and the inverse is
\[
 u=W+\one,\qquad w=U+u+(1+y\cdot W)\one.
\]
The chosen affine slice has index four in the complete
solution space: the rational $2$-torsion classes
$(\varepsilon,u,w)=(1,0,\one)$ and $(0,\one,0)$ toggle
$\varepsilon$ and $z\cdot u$, respectively.  Consequently
the solution space has $4|\ker M|$ elements, proving
\eqref{eq:Selrank}.
\end{proof}

\begin{lemma}[Bordered determinant]\label{lem:border}
Suppose $\beta\mid n$, $\beta\equiv1\pmod8$, and
$d\mid n$, $d\equiv5\pmod8$.  Let $u=\one_B$ be the
indicator of the prime factors of $\beta$ and $w$ the
indicator of those of $d$.  If $(\beta,d)\in\Sel(E_n)$,
then
\[
       k=\binom{u+w}{u}\in\ker M,\qquad b^\top k=1.
\]
If $\corank M=1$, then $\det M_{5a}=1$.
\end{lemma}
\begin{proof}
Here $y\cdot u=z\cdot u=0$, $z\cdot w=1$, and
$y\cdot w=y\cdot u=0$.  The class corresponds to
$(0,u,w)$ in \eqref{eq:S1}--\eqref{eq:S3}.  Adding the
$2$-torsion class $(0,\one,0)$ and applying the bijection
in Lemma~\ref{lem:selmer} gives the asserted kernel vector.
Also
\[
 b^\top k=(y+z)\cdot(u+w)=z\cdot w=1.
\]
If $\ker M=\langle k\rangle$, a vector $(v,t)$ in the kernel
of $M_{5a}$ satisfies $Mv+tb=0$, $b^\top v=0$.
Pair the first equation with $k$ to get $t=0$;
then $v\in\langle k\rangle$ and $b^\top v=0$ force $v=0$.
Hence $M_{5a}$ is nonsingular.
\end{proof}

\section{Why the first genus sum is the bordered determinant}\label{sec:det}

For a subset $T$ of prime indices, let $A_T$ be the
R\'edei matrix computed with only the indices in $T$:
its off-diagonal entries are inherited from $A$ and
its diagonal entries are recomputed as its row sums.
Set
\[
 G(T)=\det\bigl(A_T\text{ with one column replaced by }z_T\bigr)
       \in\FF.
\]
It is independent of the replaced column because $A_T\one=0$.
The R\'edei genus criterion identifies $G(T)=g(n_T)$
whenever $n_T\equiv1\pmod4$.  In particular this applies
to type-1 and type-5 blocks.  We record the determinant
identity used in both \cite{Smith} and \cite[\S6.2]{Jaswal},
including its combinatorial derivation.

\begin{proposition}[Row 5a, in explicit matrix form]\label{prop:smith}
For any collection of prime indices of total type 5
(i.e. $\sum y_i=0$, $\sum z_i=1$), one has
\begin{equation}\label{eq:smith}
                         \det M_{5a}=\Sigma_{5a}(n).
\end{equation}
More generally, for a configuration $T$ of total type 1,
\[
 \det M(T)=\mathcal P_1(T),\qquad
 \mathcal P_1(T):=
 \sum_{\substack{\pi\text{ partition of }T\\
                  \text{all blocks of type 1}}}\prod_{U\in\pi}G(U),
 \quad \mathcal P_1(\varnothing)=1.
\]
\end{proposition}
\begin{proof}
We reproduce the determinant recursion, to identify exactly
which part of the literature is needed.  For a symmetric
matrix in characteristic two, terms in the permutation
expansion cancel in inverse pairs unless the permutation
is an involution.  Hence its determinant is the weighted
sum of matchings (fixed vertices carry diagonal weights,
paired vertices carry the off-diagonal weight).

Apply this to $M$.  Choose $P$ as the set of indices whose
$U$-vertices are paired with another $U$-vertex and $R$ as
those whose $W$-vertices are fixed.  The $U$--$W$ pairs
form a bijection $I\setminus R\to I\setminus P$, weighted
by the corresponding minor of $A$.  By the all-minors
matrix-tree theorem \cite{Chaiken}, that minor is a sum over
rooted spanning forests; each tree carries a chosen root,
a weight $z$ at the root, and marked vertices of weight $y$.
When summed over its marks, an ordinary tree contributes
$Y_T=\sum_{i\in T}y_i$; the tree containing the distinguished
vertex $U_c$ in the principal minor contributes $1$;
and the special tree meeting the added border vertex
contributes $Y_T+Z_T$ (here $Z_T=\sum_{i\in T}z_i$).
The parity of the number of trees gives the three resulting
recursions (all sums in $\FF$):
\begin{align}
 \det M(I)&=\sum_{\substack{\pi\text{ partition of }I\\
                         \text{all blocks type 3}}}\prod_{T\in\pi}G(T)
                &&\text{if }Y_I=0;                                      \label{eq:rec1}\\
 \det M(I)^{(U_c)}
   &=\sum_{\substack{c\in T\subseteq I\\T\text{ type 1}}}
          G(T)\det M(I\setminus T)
                &&\text{if }I\text{ is type 1};                         \label{eq:rec2}\\
 \det M_{5a}(I)
   &=\sum_{\substack{T\subseteq I\\T\text{ type 5}}}
          G(T)\det M(I\setminus T)
                &&\text{if }I\text{ is type 5}.                         \label{eq:rec3}
\end{align}
In \eqref{eq:rec1} the type-3 blocks carry $Y_T=Z_T=1$;
for $Y_I=1$ the determinant is zero.  For completeness,
the all-minors formula itself follows by expanding every
diagonal entry $a_{ii}=\sum_{j\ne i}a_{ij}$: a directed
cycle is canceled by toggling that cycle in the determinant
permutation, and an acyclic map contributes exactly when
each rooted tree contains a unique marked vertex.
This is the characteristic-two form of Chaiken's identity.

For type-1 $I$, $\corank M(I)$ is even.  Indeed, with
$N=M(I)$, its diagonal vector is $\delta=(0,z)$,
and $N(\one,\one)^\top=\delta$; furthermore
$(\one,\one)N(\one,\one)^\top=Z_I=0$.
The associated symmetric bilinear form is alternating
on the orthogonal hyperplane of $(\one,\one)$, implying
even nullity.  Thus, if $N$ is singular, its adjugate
vanishes.  If $N$ is invertible, its diagonal satisfies
\(\diag(N^{-1})=N^{-1}\delta=(\one,\one)^\top\),
so $\det N^{(U_c)}=\det N$ in either case.
Equation~\eqref{eq:rec2} now provides the recursion for
$\mathcal P_1(I)$, by sorting set partitions according to
the type-1 block containing $c$.  Induction on $|I|$
yields $\det M(I)=\mathcal P_1(I)$, with the empty
configuration as base case.

For total type 5, the complement of a type-5 subset
has type 1.  Substitute the first identity into
\eqref{eq:rec3}; the right side enumerates precisely
partitions with one type-5 block and all other blocks
type 1.  R\'edei's genus criterion identifies their
weights with \eqref{eq:genus}.  This proves
\eqref{eq:smith}.
\end{proof}

\section{The five-index core and prime compression}\label{sec:core}

A configuration can be specified abstractly using the
vectors $y,z$ and off-diagonal additive symbols $a_{ij}$,
subject only to \eqref{eq:reciprocity}.  The rank and
determinant identities of the preceding sections are
polynomial identities in these bits and therefore can be
established before any primes are selected.

\begin{lemma}[Five-index core]\label{lem:five}
Let the indices be $\{\widehat\beta,r_1,r_2,r_3,r_4\}$,
with $(y,z)=(0,0)$ on $\widehat\beta$, $(1,1)$ on
$r_1,r_2,r_3$, and $(1,0)$ on $r_4$.  Require
\[
 a_{r_i r_j}=0\ (i<j),\qquad
 (a_{\widehat\beta,r_1},\dots,a_{\widehat\beta,r_4})
              =(0,0,1,1)
\]
and use reciprocity for reversed entries.  Then the
complete solutions of \eqref{eq:S1}--\eqref{eq:S3} are
\begin{equation}\label{eq:base_solution}
 u=(u_0,a,a,a,a),\qquad
 w=(\varepsilon,\varepsilon,\varepsilon,
       \varepsilon+u_0+a,\varepsilon+u_0+a),
     \quad (\varepsilon,u_0,a)\in\FF^3.
\end{equation}
Consequently $\corank M=1$ and the principal minor of $M$
obtained by deleting coordinate $U_{r_3}$ has determinant one.
\end{lemma}
\begin{proof}
The R\'edei matrix has rows, in the displayed order,
\[
 \begin{pmatrix}
 0&0&0&1&1\\0&0&0&0&0\\0&1&1&0&0\\
 1&1&1&1&0\\1&1&1&1&0
 \end{pmatrix}.
\]
Substitution into (S1) yields
\(u_3=u_4\), \(w_1=\varepsilon\),
\(w_2=\varepsilon+u_1+u_2\), and
\(w_3=w_4=\varepsilon+u_0+u_1+u_2+u_3\).
The remaining equations (S2) give
$u_1=u_2=u_3$, $w_0=\varepsilon$; (S3) is automatic.
This is exactly \eqref{eq:base_solution}.
By Lemma~\ref{lem:selmer}, $\corank M=1$.
The class with $u=e_{\widehat\beta}$ and
$w=e_{r_3}+e_{r_4}$ gives the kernel vector
$(e_{\widehat\beta}+e_{r_3}+e_{r_4},
  e_{\widehat\beta})$, whose $U_{r_3}$ coordinate is
one.  A symmetric matrix of corank one has a nonsingular
principal minor upon deleting any coordinate at which
a kernel generator is nonzero: if a vector in the kernel
of the minor is extended by zero, orthogonality to the
kernel generator forces the omitted row to vanish as well.
This proves the claim.
\end{proof}

\begin{lemma}[Prime-compression determinant]\label{lem:compress}
Let $B$ be a nonempty set of indices of total type 1,
$|B|=k$.  Add $k-1$ type-1 indices
$\Pi=\{\pi_1,\ldots,\pi_{k-1}\}$ with
\[
 a_{b\pi_i}=a_{\pi_i b}=C_{bi},\qquad
  C_{bi}=[b\in\{b_1,b_{i+1}\}],
\]
where $B=\{b_1,\ldots,b_k\}$, and set all other symbols
incident to $\Pi$ equal to zero.  Let $R$ be another
disjoint index set and put
\[
                  \varsigma_x=\sum_{b\in B}a_{xb}
                                =\sum_{b\in B}a_{bx}
                    \qquad(x\in R),
\]
where the equality follows from $Y_B=0$.
Replace $B$ by one formal type-1 index $\widehat\beta$
whose symbols to $x\in R$ are $\varsigma_x$ in either
direction.  Write $I=B\sqcup\Pi\sqcup R$,
$I'=\{\widehat\beta\}\sqcup R$.
Then for any coordinate $c$ associated to an index in $R$,
\[
                   \det M(I)^{(c)}=\det M(I')^{(c)}.
\]
\end{lemma}
\begin{proof}
The columns of $C$ form a basis of the even-weight
hyperplane of $\FF^B$, with orthogonal complement
$\langle\one_B\rangle$.  In each of the two copies of
$\FF^B$ (the $U$- and $W$-coordinates), choose a basis
beginning with $\one_B$ and completing it by a
complement $H$.  In these coordinates, the $\Pi$ variables
pair nondegenerately with the $H$ variables, since the
pairing matrix is the restriction of $C^\top$ to $H$.
The Gram matrix has block form
\[
 \begin{pmatrix}
 G_{11}&G_{12}&0\\G_{21}&G_{22}&Y\\0&Y^\top&0
 \end{pmatrix}
\]
with $Y$ invertible.  Elementary row and column
elimination gives determinant $(\det Y)^2\det G_{11}$,
and the same equality after deleting a coordinate in $R$.
Over $\FF$ the scalar $(\det Y)^2$ equals one.
The two retained $B$-vectors are the sums $U_B,W_B$.
Their self-pairings vanish because $Y_B=Z_B=0$,
and their pairings with each $x\in R$ give precisely
$\varsigma_x$ in the $U$--$W$ cross block; the diagonal
of the R\'edei matrix on $R$ receives $\varsigma_x$
from $B$.  Thus $G_{11}$ is exactly $M(I')$,
and its minor agrees as claimed.
\end{proof}

\section{Independent Legendre-symbol freedom}\label{sec:free}

The following key observation makes a single twist work
for arbitrarily many pairwise-disjoint auxiliary supports.

\begin{lemma}[One free-support determinant]\label{lem:free}
Let $I_0$ be a configuration of type 5, choose
$q_0\in I_0$, and let $J$ be a disjoint new index set
of total type 1.  Suppose
\[
 \sum_{j\in J}a_{xj}=\sum_{j\in J}a_{jx}=0
           \quad\text{for every }x\in I_0\setminus\{q_0\}.
\]
The bits $t_j=a_{j q_0}$ may vary independently,
all other data and reciprocity being fixed.
If $c$ is a coordinate belonging to $I_0\setminus\{q_0\}$
and $\det M(I_0)^{(c)}=1$, then for at least one
$t\in\FF^J$ one has
\[
                   \det M(I_0\sqcup J)^{(c)}=1.
\]
\end{lemma}
\begin{proof}
Varying $t_j$ changes $a_{q_0j}$ and $a_{jq_0}$
together by one; since the diagonal entries of $A$ are
row sums, it changes $A$ by $b_jb_j^\top$ where
$b_j=e_{q_0}+e_j$.  Change the basis in both the $U$
and $W$ copies by replacing $e_{q_0}$ with
$f=e_{q_0}+\sum_{j\in J}e_j$ and leaving all other
basis vectors fixed.  In the new coordinates the
variable $t_j$ occurs only in the symmetric positions
$(U_j,W_j)$ and $(W_j,U_j)$.

The principal submatrix on the old $I_0$ coordinates
is precisely $M(I_0)$.  Indeed, all couplings of $f$
with an old coordinate $x\ne q_0$ are unchanged because
the sums of the incident symbols vanish; the diagonal
involving $f$ is unchanged because $Y_J=Z_J=0$ and
$\sum_{j\in J}y_j=0$; the other old entries are unaffected.
The change of basis fixes the coordinate $c$.

Expand the determinant with the $c$th row and column
deleted.  It is a function $F(t)$ on $\FF^J$ and so has a
unique multilinear polynomial representation.
Its coefficient of the top monomial
$\prod_{j\in J}t_j$ is $\det M(I_0)^{(c)}$: a term
containing all the $t_j$ must match each $U_j$ with its
corresponding $W_j$, after which the residual determinant
is exactly the old principal minor.  This coefficient
is one, hence $F$ is not the zero multilinear polynomial.
A multilinear polynomial over $\FF$ is identically zero
as a function on $\FF^J$ if and only if all its
coefficients vanish (successive evaluation at 0 and 1
proves this).  Thus $F(t)=1$ for some $t$.
\end{proof}

\begin{lemma}[Disjoint simultaneous support]\label{lem:disjoint}
Fix $I_0$, $q_0$, $c$ as in Lemma~\ref{lem:free}.
For each $1\le j\le s$ let $J_j$ be a separate new
index set satisfying that lemma's conditions, with all
$J_j$ pairwise disjoint.  Then the free bits
$t_q$ for $q\in\bigcup_jJ_j$ can be selected once so that
\[
            \det M(I_0\sqcup J_j)^{(c)}=1
                              \qquad(1\le j\le s).
\]
All bits for pairs between distinct $J_i,J_j$ may be
fixed arbitrarily: they do not enter any displayed minor.
\end{lemma}
\begin{proof}
For each $j$ the determinant in question depends only
on the variables indexed by $J_j$, with a nonzero top
coefficient.  By Lemma~\ref{lem:free} choose a satisfying
value on each $J_j$.  Since the supports are disjoint,
these choices do not conflict, and the remaining symbols
are immaterial to the separate determinants.
\end{proof}

\section{Proof of the simultaneous-twist theorem}\label{sec:proof}

\begin{proof}[Proof of Theorem~\ref{thm:sim}]
Let $B=\{b_1,\dots,b_k\}$ be the set of prime factors
of $\beta$ and $J_j=\{q:q\mid\theta_j\}$.
The sets $J_j$ are disjoint and disjoint from $B$.
Each $B$ and each $J_j$ has type 1, since its product
is congruent to one modulo eight.
For $1\le i\le k-1$ put
\[
 C_{b i}=[b\in\{b_1,b_{i+1}\}],\qquad
 c_b=\sum_{i=1}^{k-1}C_{bi},\qquad
 z_b=\sgnbit{2}{b}.
\]
If $k=1$ there will be no $\pi_i$; every formula below
still makes sense with the empty sum.

\emph{Step 1: prescribe the abstract prime symbols.}
Introduce auxiliary indices $\pi_1,\ldots,\pi_{k-1}$
of residue $1\pmod8$, and $r_1,r_2,r_3$ of residue
$3\pmod8$ and $r_4$ of residue $7\pmod8$.
Set all symbols between the $\pi$ indices to zero,
and all symbols between $\pi$ and any $r_h$ or $q\in J_j$
to zero.  Set
\begin{align}
  a_{b\pi_i}=a_{\pi_i b}&=C_{bi},\label{eq:pi_bits}\\
  \bigl(a_{b r_1},a_{b r_2},a_{b r_3},a_{b r_4}\bigr)
   &=\bigl(c_b+z_b,\ 0,\ [b=b_1],\ z_b+[b=b_1]\bigr),
                                                   \label{eq:r_bits}\\
  a_{r_i r_h}&=0\qquad(i<h).\label{eq:rr_bits}
\end{align}
All reverse symbols are determined by reciprocity.
For $q\in J=\bigcup_j J_j$, prescribe
\begin{align}
 a_{q\pi_i}=a_{\pi_iq}&=0,\\
 a_{q r_h}&=t_q+y_q\quad(h=1),\qquad
                   a_{r_1q}=t_q,\label{eq:t_bits}\\
 a_{r_hq}&=0\quad(h=2,3,4),
          \qquad a_{q r_h}=y_q\quad(h=2,3,4).
\end{align}
Here the only free bits are $t_q\in\FF$.
The symbols within $B$ and within each $J_j$ are the
actual Legendre symbols of the given integers.

The pairwise admissibility hypothesis makes
\[
 \sum_{q\in J_j}a_{bq}=
 \sum_{q\in J_j}a_{qb}=0\quad(b\in B),
 \qquad \sum_{q\in J_j}y_q=
       \sum_{q\in J_j}z_q=0.
\]
The conditions of Lemma~\ref{lem:free} hold for
$I_0=B\sqcup\Pi\sqcup\{r_1,r_2,r_3,r_4\}$,
$q_0=r_1$, and each separate $J_j$: for $r_2,r_3,r_4$
the required reverse sums equal $\sum_{q\in J_j}y_q=0$.
We shall choose all the bits $t_q$ after checking
invertibility for $I_0$.

\emph{Step 2: the prescribed Selmer class exists on every support.}
Set
\[
             u=\one_B,\qquad
             w=e_{r_3}+e_{r_4},\qquad \varepsilon=0.
\]
These are the square-class coordinates for
$(\beta,d)$ with $d=r_3r_4$.
We check \eqref{eq:S1}--\eqref{eq:S3} on $I_0$
and on each $I_0\sqcup J_j$.
Because $Z_B=Y_B=0$, we have $z\cdot u=0$ and
$y\cdot u=y\cdot w=0$.
For (S1), the rows indexed by $b\in B$ give
\[
 (A\one_B)_b
 =c_b+(c_b+z_b)+0+[b=b_1]+(z_b+[b=b_1])=0;
\]
the additional $J_j$ contribution is
$[\theta_j/b]=0$.  Each row indexed by $\pi_i$
has the even column sum $\sum_bC_{bi}=0$.
By reciprocity and $Y_B=0$ the rows indexed by $r_h$
give $(A\one_B)_{r_h}=0,0,1,1$, exactly
$(D_yw)_{r_h}$.  Finally the $q\in J_j$ rows give
$[\beta/q]=0$.
For (S2), the row $b\in B$ has
$(Aw)_b=[d/b]=[r_3/b]+[r_4/b]=z_b$,
which is canceled by $(D_z u)_b=z_b$.
Rows $q\in J_j$ and $\pi_i$ have $(Aw)_q=(Aw)_{\pi_i}=0$.
For $r_3$ and $r_4$, use the row sums in the diagonal
entry of $A$ and \eqref{eq:rr_bits}; one obtains
$(Aw)_{r_3}=(Aw)_{r_4}=1$, which is canceled by
$(D_yw)_{r_3}=(D_yw)_{r_4}=1$.
The other $r_h$ rows give zero.  For the extended
supports, the only relevant extra diagonal contributions
are at $r_h$ for $h=3,4$, where
$\sum_{q\in J_j}a_{r_hq}=0$ by construction (and
$\sum_{q\in J_j}a_{q r_h}=Y_{J_j}=0$ by reciprocity).
The diagonal at $r_1$ may depend on the free bits, but
$w_{r_1}=0$, so it does not enter (S2).
Finally (S3) is $1+1=0=Y_B$.  Thus
\begin{equation}\label{eq:class_all}
            (\beta,r_3r_4)\in\Sel(E_{I_0})
               \cap\bigcap_{j=1}^s\Sel(E_{I_0\sqcup J_j})
\end{equation}
for \emph{every} choice of the free $t_q$.
Here $E_I$ temporarily means the formal Selmer
configuration with prime set $I$.

\emph{Step 3: the base principal minor is invertible.}
By Lemma~\ref{lem:compress}, the determinant of
$M(I_0)^{(U_{r_3})}$ equals the same determinant on the
five-index compressed configuration.  Indeed
\[
 \varsigma_{r_1}=\sum_b(c_b+z_b)=0,\quad
 \varsigma_{r_2}=0,\quad
 \varsigma_{r_3}=1,\quad
 \varsigma_{r_4}=1;
\]
the first equality uses the even column sums of $C$
and $Z_B=0$.  Thus the compressed data are precisely
those of Lemma~\ref{lem:five}, and
\begin{equation}\label{eq:baseminor}
                      \det M(I_0)^{(U_{r_3})}=1.
\end{equation}

\emph{Step 4: select all free bits simultaneously.}
Lemma~\ref{lem:disjoint} with $q_0=r_1$ and
$c=U_{r_3}$ now gives one choice
$t=(t_q)_{q\in J}$ for which
\begin{equation}\label{eq:allminor}
 \det M(I_0\sqcup J_j)^{(U_{r_3})}=1
                         \quad(1\le j\le s).
\end{equation}
By \eqref{eq:class_all} and Lemma~\ref{lem:border},
all matrices $M$ involved have the nonzero kernel vector
\(\binom{\one_B+w}{\one_B}\), whose $U_{r_3}$ coordinate
is one.  The nonzero principal minors
\eqref{eq:baseminor}--\eqref{eq:allminor} therefore
force $\corank M=1$ on all supports.  Hence
\(\dim\Sel=3\) by Lemma~\ref{lem:selmer}; then
Lemmas~\ref{lem:border} and Proposition~\ref{prop:smith}
give $\Sigma_{5a}=1$ on every support.

\emph{Step 5: realize the abstract configuration by primes.}
Choose the primes in the order
\[
                  \pi_1,\ldots,\pi_{k-1},
                  r_1,r_2,r_3,r_4.
\]
At the moment each prime $p$ is chosen, its residue
modulo eight and its Legendre symbols modulo every
previously selected prime and every prime of
$\beta\theta_1\cdots\theta_s$ are specified by
\eqref{eq:pi_bits}--\eqref{eq:t_bits}.
These are independent congruence conditions on $p$:
for each already fixed odd prime $q$, choose one of its
nonzero residue classes having the desired quadratic
character, and choose the required odd class mod eight.
The Chinese remainder theorem gives a reduced residue
class modulo the product, and Dirichlet's theorem provides
arbitrarily large new primes in that class.  Reciprocity
supplies the desired reverse symbols.
Set
\[
          e=\beta\pi_1\cdots\pi_{k-1}r_1r_2r_3r_4,
                  \qquad d=r_3r_4.
\]
The product is positive and squarefree, prime to all
$\theta_j$, and
\[
 e\equiv 1\cdot1^{k-1}\cdot3^3\cdot7
        \equiv5\pmod8,\qquad
 d\equiv3\cdot7\equiv5\pmod8.
\]
Thus $e\theta_j\equiv5\pmod8$ and all the preceding
finite-field calculations apply to the actual curves.
The prime factor count is
$\omega(e)=k+(k-1)+4=2k+3$.

\emph{Step 6: pass from Selmer classes to rational points.}
Apply Proposition~\ref{prop:tyz} to each of
$e,e\theta_1,\ldots,e\theta_s$.
This proves analytic and Mordell--Weil rank one,
trivial $2$-primary $\Sha$, and surjectivity of the
Kummer map.  Equation~\eqref{eq:class_all} is therefore
represented on each curve by a point of infinite order.
Because such a point is not one of the three nontrivial
$2$-torsion points, its first Kummer coordinate is the
square class of its abscissa, proving the prescribed
$x$-square class.

\emph{Step 7: quantitative infinitude.}
Fix the auxiliary primes $\pi_1,\ldots,\pi_{k-1}$,
$r_1,r_2,r_3$ selected above and their already fixed
Legendre symbols.  At the final stage, all conditions on
$r_4$ specify a reduced residue class
$a\pmod M$, where $M$ is an odd prime product times
eight; it is nonempty by the same CRT argument.
Every sufficiently large prime $r\equiv a\pmod M$
can replace $r_4$ with no change in any matrix bit.
Let $C=\beta\pi_1\cdots\pi_{k-1}r_1r_2r_3$.
Then every such $e=Cr$ has all the asserted properties.
By the prime number theorem in arithmetic progressions,
\[
 \#\{r\le X/C:r\equiv a\pmod M\}
       \sim\frac{\operatorname{Li}(X/C)}{\varphi(M)}
       \asymp_{C,M}\frac{X}{\log X}.
\]
Removing finitely many forbidden primes changes only a
bounded term.  This proves the claimed lower bound.
\end{proof}

\section{Exact checks and a concrete three-curve instance}\label{sec:checks}

The preceding proof does not depend on numerics.  The checks
below are useful for detecting transcription mistakes in the
finite-field steps; they are evidence, not replacements for
proof.

\subsection{Finite-field checks}

An independent script builds matrices $A,M,M_{5a}$ over
$\FF$, computes the rank of the Selmer system
(S1)--(S3), evaluates the first genus sum by enumerating
set partitions and R\'edei determinants, and tests
\[
  \dim\operatorname{Sol}(S1\text{--}S3)=2+\corank M,
           \qquad\det M_{5a}=\Sigma_{5a}.
\]
The five-index base configuration has Selmer dimension
three, $\corank M=1$, determinant-one $U_{r_3}$ principal
minor, and $\det M_{5a}=1$.
The two displayed identities were also checked on
independently randomized quadratic-reciprocity-compatible
configurations of type 5: $100$ trials each for
$1\le r\le4$, and $40$ trials each for $r=5,6$.
No counterexample occurred.  The standalone source file \texttt{check\_row5a.py} accompanies this
note and verifies all the claims displayed in this subsection.
It additionally runs $420$ random type-1 matrix/partition checks.
All arithmetic is exact modulo two; primality and Legendre symbols
for the numerical example are verified with \texttt{sympy}.

\subsection{A simultaneous instance with distinct auxiliary twists}

Take
\[
        \beta=17,\qquad
        \theta_1=89,\qquad\theta_2=137.
\]
All three are primes $\equiv1\pmod8$, and
$(17/89)=(17/137)=+1$, so the pairs are admissible.
Here $B$ has one element and there are no $\pi_i$.
The choice of free bits $t_{89}=t_{137}=1$ is realized by
\[
       (r_1,r_2,r_3,r_4)=(43,443,107,4519),
\]
all prime, with respective residues $3,3,3,7\pmod8$.
It gives
\[
 \boxed{\ e=17\cdot43\cdot443\cdot107\cdot4519
          =156583941989,\qquad d=107\cdot4519=483533.\ }
\]
Exact F$_2$ arithmetic gives the following verifiable
rank/determinant certificates:
\begin{center}
\begin{tabular}{@{}lccc@{}}
\toprule
$n$ & $\dim_{\FF}\Sel(E_n)$ & $\det M_{5a}(n)$
                                  & $\Sigma_{5a}(n)$\\
\midrule
$e$      & $3$ & $1$ & $1$\\
$89e$    & $3$ & $1$ & $1$\\
$137e$   & $3$ & $1$ & $1$\\
\bottomrule
\end{tabular}
\end{center}
For each of these three curves the rank-one conclusion,
triviality of $\Sha[2^\infty]$, and realization of
$(17,483533)$ now follow from the proved analytic and
Kummer implications.  The check has not computed the
individual infinite-order points or their heights.

\section{Validity assessment and implications for Diophantine complexity}\label{sec:audit}

\textbf{Theorem 1.2 / Section 6.}  The argument has four
nontrivial links: the exact local $2$-Selmer system;
the bordered-determinant/first-genus-sum identity;
the five-index compression and nonzero free-bit minor;
and the published Tian--Yuan--Zhang criterion.
The first three are checked by the proofs above and by
independent exact finite-field tests; the fourth was
compared with \cite[Theorem~1.4]{TYZ} in the published
journal text.  No essential defect was found.  In
particular the theorem does not need the unpublished
Proposition~N of the row-5b route, an unpublished
restricted $2$-converse, or a Birch--Swinnerton-Dyer
conjecture.  It \emph{does} use the published deep
Tian--Yuan--Zhang theorem and Gross--Zagier.

\textbf{The rest of the paper.}  Sections~2--5 and~7--8
are predominantly a reconstruction of earlier project
methods: the rank-one curve over $\QQ(\sqrt2)$, a pole
predicate, formal logarithm comparisons, six-contact
heights, the eighteen-prime sieve/prime-pattern argument,
and elliptic Lucas growth.  The switch to the simpler
congruent-number pole equations (without multiplier lists),
a different modularity proof via Serre's theorem rather than
Pan's, the use of three slopes to make the finite coefficient
set uniform in the input, and the scalar-condition guard are
useful technical adaptations.  They are not identified here
as new quantitative Diophantine theorems.  This audit of their statements, dependency
structure, and key arithmetic interfaces found no
specific counterexample or decisive contradiction;
it is \emph{not} an independent line-by-line certification
of every global geometric and analytic proof in those
sections.  The proof of Theorem~\ref{thm:sim} needs none
of those sections.

\textbf{One nonfatal error in Section~7.}  The CM case of
Proposition~7.3 (page~20 in the preprint) asserts that
finitely many possible CM discriminants imply finitely many
CM elliptic curves \emph{over $\QQ$}.  This is false:
for example, the curves $y^2=x^3-m^2x$, as $m$ ranges
over squarefree positive integers, give infinitely many
$\QQ$-isomorphism classes with $j=1728$.  The intended
height conclusion is repairable: a fixed CM discriminant
permits only finitely many $\overline{\QQ}$-isomorphism
classes (finitely many $j$-invariants), and the \emph{stable}
Faltings height is invariant under geometric isomorphism.
The degree-bounded isogeny step then proceeds from the
height bound, not from finiteness of rational twists.
This correction is not a counterexample to Theorem~1.3.

\textbf{Depth limitation.}  The long argument in Section~7
imports several difficult results on quaternionic Shimura
curves, descent of compatible Galois representations, Serre
modularity and Faltings heights.  Its dependency chain and
central quantitative exponents were examined, but the present
audit does not certify each of those imports or all their
normalizations independently.  The same caveat applies to
the fully uniform prime-pattern asymptotics in Section~8.
Accordingly, the absence of a found fatal flaw in the
full theorem is not a claim of complete external verification.

\textbf{What this does and does not optimize.}  A common
$e$ across several compatible pole systems creates a
potential opportunity for sharing witness parameters.
However the present theorem gives no elimination of a
retained rational variable, no degree bound for a defining
polynomial, and no bound on the heights of the rational
points.  It would be incorrect to announce an improved
Diophantine variable count from this theorem alone.
A future application must explicitly construct the
new existential formula, check its zero set in both
directions, and count all retained coordinates.

\textbf{Outstanding extension.}  When the $\theta_j$
share prime factors, their free-bit subsets overlap.
The individual determinant polynomials may have disjoint
sets of satisfying assignments, so nonvanishing of their
top coefficients does not on its own show that all can
be made nonzero simultaneously.  The theorem above
therefore makes no claim in this regime.  Removing this
condition would require a new compatibility argument or
new auxiliary-prime slots.

\begin{thebibliography}{99}
\bibitem{Jaswal}
M.~Jaswal, \emph{A proof that $\ZZ$ is Diophantine in $\QQ$
through congruent-number twists}, hexagon:2611.00331v2,
12 November 2026,
\url{https://hexagonmath.org/pdf/2611.00331v2}.
\bibitem{TYZ}
Y.~Tian, X.~Yuan and S.-W.~Zhang,
\emph{Genus periods, genus points and congruent number problem},
Asian Journal of Mathematics \textbf{21} (2017), no.~4,
721--774, especially Theorem~1.4 and Theorem~3.5,
\url{https://archive.ymsc.tsinghua.edu.cn/pacm_download/355/9144-cn.pdf},
DOI: \href{https://doi.org/10.4310/AJM.2017.v21.n4.a5}{10.4310/AJM.2017.v21.n4.a5}.
\bibitem{Smith}
A.~Smith, \emph{The congruent numbers have positive natural density},
preprint arXiv:1603.08479, Theorem~2.2 and Table~2
(row 5a); the determinant argument used here is also
reproduced in \cite[\S6.2]{Jaswal}.
\bibitem{Chaiken}
S.~Chaiken, \emph{A combinatorial proof of the all minors
matrix tree theorem}, SIAM Journal on Algebraic and
Discrete Methods \textbf{3} (1982), 319--329.
\bibitem{Redei}
L.~R\'edei, \emph{Arithmetischer Beweis des Satzes \"uber die
Anzahl der durch vier teilbaren Invarianten der absoluten Klassengruppe
im quadratischen Zahlk\"orper}, Journal f\"ur die reine und
angewandte Mathematik \textbf{171} (1934), 55--60,
\url{https://doi.org/10.1515/crll.1934.171.55}; for the precise
matrix form see \cite[Proposition~6.7]{Jaswal}.
\bibitem{GZ}
B.~Gross and D.~Zagier, \emph{Heegner points and derivatives
of $L$-series}, Inventiones Mathematicae \textbf{84}
(1986), 225--320.
\bibitem{Glazer}
E.~Glazer and collaborators,
\emph{Diophantine definability of $\ZZ$ in $\QQ$},
project manuscript and verified research notes,
\url{https://github.com/elliotglazer/diophantine-z-q},
particularly \texttt{docs/research/tyz-5b.md} and
\texttt{blueprint/roadmap/rank/literature/genus-sum-5a.md}.
\bibitem{Brox}
J.~Brox, \emph{Diophantine $\ZZ$ in $\QQ$: improvements
and quantitative constructions}, project repository,
\url{https://github.com/JoseBrox/diophantine-z-q_improvements}.
\bibitem{Dirichlet}
Dirichlet's theorem and the prime number theorem for
arithmetic progressions, used in their standard forms.
\end{thebibliography}
\end{document}
